\section{Modern Pseudorandom Number Generators}

While LFSRs and LFGs provide foundational algebraic methods for generating pseudorandom sequences, modern applications require generators that execute with minimal latency in software while exhibiting strong empirical statistical properties. This chapter explores two highly efficient families of generators: \textbf{Xorshift} (and its scrambled variants) and \textbf{PCG} (Permuted Congruential Generator).

\subsection{The Xorshift}

Introduced by George Marsaglia in 2003, the Xorshift generator operates exclusively via hardware-efficient logical shift and bitwise XOR operations. Mathematically, it implements a linear transformation over the vector space $\mathbb{F}_2^w$, making it a software-efficient analogue of the LFSR.

\subsubsection{The Recurrence Relation}


Given a non-zero $w$-bit integer state $\mathbf{x}_n$, the subsequent state is generated by applying three successive shift-and-XOR operations using integer constants $a$, $b$, and $c$:
\begin{align*}
    \mathbf{y}_1 &= \mathbf{x}_0 \oplus (\mathbf{x}_0 \ll a) \\
    \mathbf{y}_2 &= \mathbf{y}_1 \oplus (\mathbf{y}_1 \gg b) \\
    \mathbf{x}_1 &= \mathbf{y}_2 \oplus (\mathbf{y}_2 \ll c)
\end{align*}


\subsubsection{Matrix Representation and Periodicity}
Rigorously, a $w$-bit word is a vector in $\mathbb{F}_2^w$. A left shift by $a$ positions equates to multiplying the state vector by a nilpotent $w \times w$ matrix $L^a$, where $L$ possesses ones on the first subdiagonal. Similarly, a right shift by $b$ is multiplication by $R^b = (L^T)^b$. Since XOR acts as standard vector addition over $\mathbb{F}_2$, the global state transition is governed by the linear map $\mathbf{x}_{n+1} = M\mathbf{x}_n$, where the transformation matrix $M$ is:
$$
M = (I + L^c)(I + R^b)(I + L^a)
$$

\begin{theorem}
An Xorshift generator defined by the state transition matrix $M$ over $\mathbb{F}_2^w$ has a maximal cycle length (period) of $2^w - 1$ if and only if $M$ has multiplicative order $2^w - 1$, equivalently if its minimal polynomial is primitive of degree $w$.
\end{theorem}
When the minimal polynomial is primitive, the state space partitions into exactly two cycles: a trivial fixed point at the zero vector (length 1), and a single maximum-length cycle containing all $2^w - 1$ non-zero states.

\subsubsection{Variants of the Xorshift}

Despite achieving maximum periodicity, pure linear Xorshift generators fail rigorous empirical statistical tests (e.g., \textit{MatrixRank} in TestU01). Because the transition function is linear over $\mathbb{F}_2$, consecutive outputs exhibit deterministically bounded matrix ranks and low-weight artifacts. To resolve this, "scrambled" variants append a non-linear output permutation to the linear state transition.


\begin{enumerate}
    \item \textbf{Xorshift*:} Appends a modular multiplication scrambler, $y_n \equiv (\mathbf{x}_{n+1} \cdot M_{\ast}) \pmod{2^w}$, to the linear state transition. This introduces non-linear bit significantly improving empirical statistical quality.

    \item \textbf{Xorshift+:} Expands the state space to multiple words and computes the output via modular addition: $y_n \equiv (\mathbf{s}_0 + \mathbf{s}_1) \pmod{2^w}$. However, because the initial carry-in is zero, the least significant bit remains strictly linear ($s_{0,0} \oplus s_{1,0}$), causing it to fail bit-reversed statistical tests.

    \item \textbf{Xoshiro (xor, shift, rotate) / Xoroshiro (xor, rotate, shift, rotate):} Integrates circular rotations into transition for faster bit diffusion across 128-bit or 256-bit states. It eliminates the lower-bit linearity flaw of Xorshift+ by utilizing dual-domain scramblers.
\end{enumerate}


\subsection{The PCG}

Introduced by Melissa E. O'Neill in 2014, the PCG framework operates exclusively via hardware-efficient modular arithmetic and bitwise permutation operations. Mathematically, it implements a non-linear transformation over the ring $\mathbb{Z}/2^k\mathbb{Z}$, making it a empirically stronger alternative to classical LCG.


\subsubsection{Representation and Periodicity}

The core engine of PCG is an affine recurrence relation operating over the integer ring $\mathbb{Z}/2^k\mathbb{Z}$:
$$
S_{n+1} \equiv (a S_n + c) \pmod{2^k}
$$
By choosing a multiplier $a \equiv 1 \pmod 4$ and an odd increment $c$, the generator guarantees an unbroken maximal cycle of length $2^k$.

Crucially, the increment $c$ uniquely defines the \textit{stream} of the generator. Because any odd $c$ yields a full period, a $k$-bit state space allows for exactly $2^{k-1}$ distinct streams. This cycle topology permits developers to assign distinct streams (via different constants $c$) to parallel threads, mathematically guaranteeing distinct state trajectories for different streams.


\subsubsection{Popular variants of the PCG Framework}

Despite achieving maximum periodicity, pure linear congruential generators fail rigorous empirical statistical tests (e.g., \textit{Spectral Test} in TestU01). To resolve this issue, "permuted" variants append a non-linear output permutation to the affine state transition.

\begin{enumerate}
    \item \textbf{PCG-XSH-RR:} Uses a random rotation scrambler, $u_n = \text{ROTR}(X, r)$. This introduces non-linear bit diffusion significantly improving empirical statistical quality.

    \item \textbf{PCG-XSH-RS:} Reduces the permutation operations to a simple random shift and computes the output via logical right shifts: $u_n = X \gg r$. However, because it avoids circular rotation, the lowest bits remain partially correlated to the state, causing it to fail some bit-reversed statistical tests.

    \item \textbf{PCG-RXS-M-XS ("random xorshift, multiply, xorshift"):} Integrates an internal 64-bit multiplication into the permutation for faster bit diffusion across 32-bit or 64-bit states. 
\end{enumerate}